\sporkSection{Language Overview}{sec:LanguageOverview}

\sporkName{} extends the Python-embedded DSL design of the original Exo language.
We scope the language to procedures (programs) targetting a single CPU core driving a single CUDA device of the Ampere (\texttt{sm\_80}) or Hopper (\texttt{sm\_90a}) architecture.
The \sporkName{} compiler generates CUDA C++ that may compile to other architectures, but we have only added accelerator instructions for these architectures.
We added a \texttt{with~CudaDeviceFunction: $\sporks$} statement that distinguished CPU code (outside) from CUDA-accelerated code (the $\sporks:\sporkStmt$ inside),
and generalized the original loops into three kinds, distinguished by $\sporkloopMode$: \\
(1) \texttt{seq}$(\sporkexpr_\mathrm{lo},\sporkexpr_\mathrm{hi})$ is a sequential loop that iterates $i$ from $\sporkexpr_\mathrm{lo}$ to $\sporkexpr_\mathrm{hi}-1$.\\
(2) \texttt{cuda\_tasks}$(\sporkexpr_\mathrm{lo},\sporkexpr_\mathrm{hi})$ distributes iterations across CUDA clusters; on pre-\texttt{sm\_90} (H100) GPUs a cluster coincides with a single CTA (thread block).\\
(3) \texttt{cuda\_threads}$(\sporkexpr_\mathrm{lo},\sporkexpr_\mathrm{hi},\; \texttt{unit}=\sporkcollUnit)$ assigns iterations to subsets of threads inside the current cluster. %
The \texttt{unit} parameter is a \emph{collective unit} that specifies the shape of the participating thread set but not concrete indices—for example, $\sporkcollUnit=\,$\texttt{cuda\_warp} denotes contiguous 32-thread groups with thread IDs \(0\text{--}31\), \(32\text{--}63\), \(64\text{--}95\), etc.\ (inclusive), while $\sporkcollUnit=\,$\texttt{cuda\_thread} denotes a single thread (see Figure~\sporkTodoLink{}).

\begin{figure}[h]
\sporkCodeRule
\centering
\begin{varwidth}{\linewidth}
\sporkListing{overview_scope}
\end{varwidth}
\caption{Program scopes created by nested parallel loops.}
\label{fig:ScopeExample}
\sporkCodeRule
\end{figure}


Each loop body, conditional branch, \texttt{with}-block, or procedure defines a \emph{program scope}.
As part of \emph{collective analysis} (Section~\sporkTodoLink{}), we deduce a per-scope \emph{thread mapping} $\sporktidMap: \sporkSigma \to \mathcal{P}(\sporkTid)$.
This describes a mapping between the control variable values and the \emph{thread collective} (set of threads) assigned to execute statement instances in the scope (i.e. which threads execute which loop iterations).
We say that a scope is a $\sporkcollUnit$-scope when $\sporkcollUnit$ accurately describes the thread sets produced by the mapping.
In Figure~\ref{fig:ScopeExample}, the outermost parallel loop \texttt{cuda\_tasks(0,37)} creates \texttt{scope-1} (which is a cluster scope, and, assuming \texttt{clusterDim=1}, also a CTA scope); nesting \texttt{cuda\_threads(..., unit=\texttt{cuda\_warp})} creates \texttt{scope-2} (warp scope); and nesting \texttt{cuda\_threads(..., unit=\texttt{cuda\_thread})} creates \texttt{scope-3} (single-thread scope).
Statement instances in \texttt{scope-3} are executed by one thread.

We inherit from Exo the requirement of \emph{control-data separation}.
Variables of boolean, size, or index (loop iterator) type are of $\sporkcontrol$ type.
All others (including the contents of arrays) are of a data type (e.g. \texttt{bf16}).
This restriction stems from Exo's rewrites,
and the real synchronization checking implementation relies on this,
but the synchronization environment semantics (\sporkTodoLink{}) do not depend on this restriction.

All statements and expressions are given deterministic sequential small-step semantics (Section~\ref{sec:Semantics}).
The semantics for parallel and asynchronous constructs in \sporkName{} are similar to the \emph{serial projection} in Cilk
\cite{Cilk5} \cite{OpenCilk:serial-projection}
in that their semantics are defined identically to a similar sequential construct.
For example, the step and done rules for parallel loops and sequential loops are the same
(Figure~\ref{fig:ControlSemantics}).

Rules define steps of the form $\langle \sporks_1, \sporksigma_1, \sporktr^*_1 \rangle \to \langle \sporks_2, \sporksigma_2, \sporktr^*_2 \rangle$; the 3-tuples consist of a \emph{statement}, a \emph{value environment} giving the numeric values of array elements, and a \emph{trace} consisting of a sequence of memory access (read/mutate) events and synchronization events ascribed to certain threads.  % Forget about observers/qtl until next full section.

Continuing the sequential-or-parallel loop example, the computed value environment $\sporksigma$ is invariant of the $\sporkloopMode$.
The trace alone ``remembers'' the parallelism (indirectly, due to parallel loops influencing the deduced $\sporktidMap$), and is analyzed for correct synchronization.
Specifically, traces are constructed using the \emph{checked trace concatenation operator} $\sporkConcatChk$ (Section~\ref{sec:TraceConcatChk}), which causes the stepping to get stuck upon attempting to append a trace event that risks violating sequential-parallel equivalence.
In spite of the determinism, we call the full semantics presented \emph{parallel semantics}.
For the purposes of reasoning about rewrites, we define \emph{sequential semantics}, which are identical except for the omission of the trace and associated risk of failed trace concatenation.

% This is kind of presented without justification for now.
We say that a procedure (with statement body $\sporks$) satisfies sequential-parallel equivalence for a given input environment $\sporksigma$ when its \emph{sequential interpretation} -- $\langle \sporks, \sporksigma \rangle$ fully stepped under sequential semantics -- and its \emph{parallel interpretation} -- $\langle \sporks, \sporksigma, () \rangle$ fully stepped under parallel semantics -- either both get stuck or both step to completion (if neither get stuck, the output $\sporksigma'$ will trivially be identical, and the parallel interpretation recorded no equivalence-violating trace events).

\sporkSubsection{System Overview}{sec:SystemOverview}

% David Zhao Akeley 2026-11-04: I'm done saying "object code"
\sporkName{} programmers can either directly write optimized code or apply a sequence of scheduling transformations to simple initial code.
Both approaches grant full control over performance-critical details.
However, the scheduling approach offers significant practical advantages for performance engineering:
optimization strategies become reusable functions across similar kernels, complex transformations decompose into sequences of verified primitives that enable rapid iteration with correctness guarantees, and the compiler automatically handles complex index calculations for tiling strategies---especially valuable for GPU optimization.
Below, we describe \sporkName{}'s scheduling-based optimization process.

The scheduling flow in \sporkName{} begins with a simple initial procedure $\sporkp_0$.
% WARNING: we use $R$ ad-hoc here, not to be confused with \sporkR, which is the real numbers.
The programmer issues a series of \emph{rewrite operations} $R_1, \ldots, R_n$ to create transformed procedures $\sporkp_1, \ldots, \sporkp_n$ (Figure~\ref{fig:RewriteFlow}).

% Rewrite flow with the chain of equivalence above it and the backend checks expanded below.
% Adapted from the InstrChecksRewrite frame of 26.tex (fig_exoRewriteFlow_body.tex plus annotations);
% deliberately a separate copy, since the paper version differs enough.
\begin{figure}[h]
\sporkForceStacktrue
\sporkCodeRule
\centering
\begin{sporkfigure}
  \tikzset{
    flow node/.style={rectangle, draw, minimum height=10mm, align=center, inner sep=1.5mm},
    proc node/.style={flow node, minimum width=12mm},
    check note/.style={rectangle, draw=black!45, fill=black!3, rounded corners=1mm,
      inner sep=1.5mm, align=center, anchor=north},
    rewrite arrow/.style={spork arrow, draw=sporkBlueFg},
    rewrite label/.style={above=0.3mm, text=sporkBlueFg},
  }

  % Rewrites row.
  \node (proc0) [proc node] {$\sporkp_0$};
  \node (proc1) [proc node, right=10mm of proc0] {$\sporkp_1$};
  \node (procM) [proc node, right=20mm of proc1] {$\sporkp_{n-1}$};
  \node (procN) [proc node, right=10mm of procM] {$\sporkp_n$};
  \node (checks) [flow node, spork yellow style, text width=16mm, right=5mm of procN] {Backend\\checks};
  \node (code) [flow node, minimum width=14mm, right=5mm of checks] {C/CUDA};
  \node (rewritesLabel) [anchor=east] at ([xshift=-2mm]proc0.west) {\textbf{Rewrites:}};

  \draw [rewrite arrow] (proc0) -- node [rewrite label] {$R_1$} (proc1);
  \draw [rewrite arrow, dash pattern=on 2.5pt off 2pt] (proc1) --
    node (rewritesStar) [above=0.3mm, font=\fontsize{16}{19}\selectfont, text=sporkBlueFg] {$*$} (procM);
  \draw [rewrite arrow] (procM) -- node [rewrite label] {$R_n$} (procN);
  \path (proc0.south) -- node (rewrites) [below=0mm, text=sporkBlueFg] {Exo Rewrites} (procN.south);
  \draw [spork arrow] (procN) -- (checks);
  \draw [spork arrow] (checks) -- (code);

  % Equivalence row, above the rewrites row.
  \coordinate (eqY) at ([yshift=4mm]proc0.north);
  \node (seq0) [anchor=base] at (eqY -| proc0) {$\sporkseqInterp{\sporkp_0}$};
  \node (seq1) [anchor=base] at (eqY -| proc1) {$\sporkseqInterp{\sporkp_1}$};
  \node (seqM) [anchor=base] at (eqY -| procM) {$\sporkseqInterp{\sporkp_{n-1}}$};
  \node (seqN) [anchor=base] at (eqY -| procN) {$\sporkseqInterp{\sporkp_n}$};
  \node (parN) [anchor=base] at (eqY -| checks) {$\sporkparInterp{\sporkp_n}$};
  \node (eqLabel) [anchor=base east] at (eqY -| rewritesLabel.east) {\textbf{Equivalence:}};
  \path (seq0.base east) -- node [anchor=base] {$\equiv$} (seq1.base west);
  \path (seq1.base east) -- node [anchor=base] {$\equiv \cdots \equiv$} (seqM.base west);
  \path (seqM.base east) -- node [anchor=base] {$\equiv$} (seqN.base west);
  \path (seqN.base east) -- node [anchor=base] {$\equiv$} (parN.base west);

  % Backend checks, expanded below; the row is right-aligned with the C/CUDA box
  % and extends left under the rewrites. The cross bar sits just under "Exo Rewrites".
  \coordinate (bracketY) at ([yshift=-2mm]rewrites.south);
  \coordinate (rowTop) at ([yshift=-3.5mm]bracketY);
  \node (syncNote) [check note, text width=12mm, anchor=north east] at (rowTop -| code.east)
    {Sync\\Check};
  \node (threadsNote) [check note, text width=24mm, anchor=north east]
    at ([xshift=-4mm]syncNote.north west) {Instr Thread\\Checks};
  \node (layoutNote) [check note, text width=30mm, anchor=north east]
    at ([xshift=-4mm]threadsNote.north west) {Instr Data\\Layout Checks};
  \node (collNote) [check note, text width=32mm, anchor=north east]
    at ([xshift=-4mm]layoutNote.north west) {Collective Analysis\\($\sporktidMap$)};
  \draw [spork line, draw=sporkYellowFg]
    ([yshift=-2.5mm]bracketY -| collNote.west) -- (bracketY -| collNote.west)
    -- (bracketY -| syncNote.east) -- ([yshift=-2.5mm]bracketY -| syncNote.east)
    (bracketY -| checks.south) -- (checks.south);
  \fill [sporkYellowBg]
    ([xshift=1mm, yshift=-0.8mm]bracketY -| collNote.west)
    rectangle ([xshift=-1mm, yshift=-2mm]bracketY -| syncNote.east);

  % Informal definition of $\equiv$, in the dead space left of the backend checks.
  \node (equivNote) [check note, fill=none, align=left, font=\sporkFigSmall, anchor=north east]
    at ([xshift=-3mm]bracketY -| collNote.west)
    {$\equiv$ (functional equiv.)\\means same output $\sporksigma'$\\given same input $\sporksigma$};
\end{sporkfigure}
\caption{\sporkName{} rewrite flow and chain of equivalence.}
\label{fig:RewriteFlow}
\sporkCodeRule
\end{figure}


Exo's scheduling rewrites verify functional equivalence of $\sporkseqInterp{\sporkp_0}, \ldots, \sporkseqInterp{\sporkp_n}$ under sequential semantics (denoted $\sporkseqInterp{\cdot}$).
\sporkName{} adds new scheduling operations for parallel and synchronization constructs, but these have no impact in sequential semantics and trivially pass $\sporkseqInterp{\cdot}$-equivalence checks.
The backend checks apply only to the final procedure $\sporkp_n$, prior to translating to C and CUDA C++.
Most notably, synchronization checking ensures $\sporkseqInterp{\sporkp_n} \equiv \sporkparInterp{\sporkp_n}$, where $\sporkparInterp{\cdot}$ denotes parallel semantics.
Transitively, this completes the \emph{chain of equivalence} from $\sporkseqInterp{\sporkp_0} \equiv \cdots \equiv \sporkseqInterp{\sporkp_n} \equiv \sporkparInterp{\sporkp_n}$, guaranteeing that $\sporkp_0$'s behavior is preserved through to the parallel code generated from $\sporkp_n$.

End-to-end functional equivalence relies on these categories of reasoning:

\sporkMainKey{Exo Rewrites:}
Exo's existing rewrites implement dependency analysis-based equivalence checking, which guards against \emph{dataflow logic bugs}.
For example, \texttt{reorder\_stmts}, which reorders two statements $s_1;s_2$, is only permitted when the effects of the statements commute.

\sporkMainKey{Instruction Substitution:}
Exo's \texttt{replace} rewrite lets programmers substitute a code block with an accelerator instruction call.
\sporkName{} models instructions via their \emph{behavior} (sequential semantics), per-parameter \emph{memory type} (physical memory and layout), expected collective unit, and expected device (CPU or CUDA).
%
Exo's unification verifies that the replaced code block's behavior matches the instruction,
while the backend checks (Figure~\ref{fig:RewriteFlow}) that check instr data layout and threads ensure the instruction successfully accesses data and executes with the correct arrangement and kind of threads.
% Wording is awkward but I need to move on for now.
The two backend checks (from Exo and \sporkName{}, respectively, Section~\ref{sec:StaticInstrChecks}) and unification together prevent \emph{instruction logic errors}.

\sporkMainKey{Collective Analysis:}
This backend check ensures the compiled $\sporkp_n$ procedure's parallel loops are actually implementable (e.g. no warp-loops at thread scope) and deduces the per-scope $\sporktidMap$ (Section~\ref{sec:CollAnalysis}).
This sets the stage for the instruction thread checks and synchronization check.

\sporkMainKey{Synchronization Check:} Since parallel loops and synchronization statements are ignored during rewrites, synchronization checking must ensure that their usage is valid and does not alter observed procedure outputs or introduce race conditions.
Although we present more intuitive trace-based parallel semantics,
the real implemented synchronization checker interprets the checked procedure with equivalent (but more efficient) \emph{synchronization environment semantics} (Section~\ref{sec:SyncEnv}),
which protects the checked procedure against both \emph{inter- and intra-thread synchronization bugs} for selected concrete problem sizes.



